% !TEX program = xelatex
\documentclass{misisarticle}
\usepackage{balance}

\grnti{28.23.27}
\asjc{1702, 2207}
\oecd{1.2, 2.2}

\title{Распределённое назначение задач в гетерогенном рое БПЛА при ограниченной и переключающейся связи}

\misisauthor{Д.~В.~Басанский}{кафедра инженерной кибернетики НИТУ «МИСИС»}{Москва, Россия}{m2305618@edu.misis.ru}

\setabstract{Рассматривается распределённое назначение пространственно разнесённых задач гетерогенному рою беспилотных летательных аппаратов при ограниченной дальности связи, потерях пакетов и изменении связности сети. Цель исследования состоит в сокращении затрат на восстановление согласованного распределения после соединения ранее изолированных групп при сохранении допустимости назначений. Исследуется событийный вариант алгоритма Consensus-Based Bundle Algorithm с частичным перепланированием: возможности аппарата задают допустимость выполнения задачи, а восстановление связи запускает ограниченный добровольный пересмотр назначений у концов восстановленной линии. Обязательное устранение конфликтов выполняется независимо от этого ограничения. Представлены постановка задачи, правило пересмотра пакетов и протокол будущей проверки на синтетических сценариях. Планируется сравнение с CBBA, вариантом CBBA с частичным перепланированием и Contract Net по согласованности назначений, качеству планов и затратам обмена. Проверка доли фактически выполненных задач выделена в последующий этап моделирования движения.}

\setkeywords{гетерогенный рой БПЛА; распределение задач; мультиагентная система; ограниченная связь; переключающаяся топология; CBBA.}

\begin{document}
\maketitle

\section{Введение}

Группа беспилотных летательных аппаратов (БПЛА) может выполнять распределённые по территории операции наблюдения, обследования и доставки. Аппараты различаются оборудованием, скоростью и доступным ресурсом движения. Поэтому при назначении исполнителя требуется учитывать как стоимость достижения задачи, так и возможность её выполнения данным аппаратом.

При децентрализованном управлении агенты формируют планы на основе собственных данных и сообщений соседей. Ограниченная дальность радиосвязи и потери пакетов приводят к тому, что сведения о назначениях распространяются не одновременно. Если две группы временно не связаны, они могут независимо выбрать исполнителей одной задачи. После восстановления связи требуется согласовать эти планы, сохранив по возможности полезную часть уже выполненных расчётов.

Алгоритм Consensus-Based Bundle Algorithm (CBBA) сочетает локальное формирование пакета задач и согласование победителей распределённого аукциона~\cite{choi2009}. Частичное перепланирование исследовано в CBBA-PR~\cite{buckman2019}. В настоящей работе рассматривается применение такого механизма к восстановлению связи между группами с неодинаковыми возможностями аппаратов. Проверяемая гипотеза состоит в том, что запуск дополнительного пересмотра у ограниченного числа агентов позволяет уменьшить время восстановления согласованности и объём обмена до её достижения, сохранив приемлемое качество распределения.

Представленная версия статьи относится к этапу постановки исследования. Сформулированы модель, исследуемое правило и план экспериментов. Преимущество метода и границы его применимости предстоит установить численно.

\section{Анализ существующих работ}

\subsection{Распределённые аукционы и перепланирование}

В CBBA каждый агент последовательно добавляет задачи в пакет по предельной полезности, а затем обменивается сведениями о победителях с соседями~\cite{choi2009}. Пакет хранит порядок формирования ставок; маршрут определяет порядок посещения задач. Эти порядки могут различаться. При потере задачи агент освобождает и последующие элементы пакета, поскольку их ставки рассчитаны с учётом ранее выбранных задач. Поэтому локальное изменение назначения может вызвать дальнейший пересмотр планов.

В CBBA-PR рассматривается поступление новых задач во время работы распределённого алгоритма~\cite{buckman2019}. Пересмотр части последних элементов пакета позволяет регулировать объём перепланирования; отдельно исследуется сокращение участвующей команды. Следовательно, частичный сброс и ограничение числа участников сами по себе не являются новым результатом настоящей работы. Здесь исследуется конкретное правило запуска по восстановлению связи и его поведение при разделении гетерогенной группы.

Raja, Habibi и How рассматривают согласование назначений при ограниченных коммуникационных ресурсах, в том числе с учётом доступа к каналу и коллизий передач~\cite{raja2022}. Это показывает необходимость учитывать сетевой обмен при оценке алгоритма. В предлагаемой начальной модели потери задаются вероятностно; физический механизм коллизий и управление доступом к радиоканалу не воспроизводятся.

В качестве дополнительного метода сравнения выбран Contract Net~\cite{smith1980}. В нём назначение организуется через объявления задач, предложения исполнителей и выдачу контрактов. Роли распределяются между агентами. Для сопоставления требуется определить конкретную реализацию протокола и её поведение при недоступности агента, объявившего задачу.

Выбор функции полезности и сетевые предпосылки существенны для анализа CBBA. Современная работа Kha и соавторов уточняет условия применимости аргументов о сходимости и качестве решения~\cite{kha2025}. Поэтому гарантии исходных постановок не переносятся в данной статье на выбранную маршрутную полезность и событийный пересмотр без отдельного доказательства.

\subsection{Смежные исследования}

Киселев и Алтунян сравнивают методы обучения с подкреплением для локального планирования движения в автомобильной среде~\cite{kiselev2025}. Бахвалов и Старцев исследуют обнаружение БПЛА разных классов средствами компьютерного зрения~\cite{bakhvalov2025}. Эти работы относятся к получению информации и выполнению движения, которые в многоагентной системе предшествуют назначению задач или следуют за ним. Они не устанавливают свойства распределённого аукциона. В настоящей модели координаты задач считаются известными, а распознавание объектов и обход препятствий вынесены за её пределы.

В R-LGP граф достижимости применяется для проверки геометрической осуществимости последовательностей действий мобильного манипулятора~\cite{ly2024}. Для рассматриваемой задачи полезен общий принцип проверки допустимости до выбора назначения. Однако ограничения возможностей БПЛА и ресурса маршрута задаются здесь самостоятельно; метод R-LGP не используется как алгоритм управления роем.

\section{Постановка задачи}

\subsection{Аппараты, задачи и ограничения}

Рассматриваются множество аппаратов $\mathcal A=\{1,\ldots,n\}$ и множество задач $\mathcal J=\{1,\ldots,m\}$. На одном интервале планирования заданы исходные положения $s_i\in\mathbb R^2$, предельные длины маршрутов $E_i$ и максимальные размеры пакетов $L_i$. Высота полёта фиксирована. Модель не включает динамику приводов, препятствия и предотвращение столкновений.

Для задачи $j$ известны положение $r_j$, ценность $w_j>0$ и требования $a_j\in\{0,1\}^{K}$. Возможности аппарата задаются вектором $c_i\in\{0,1\}^{K}$. Его компонентами могут быть наличие камеры, тепловизора или средства доставки. Скорости аппаратов и длительности обслуживания понадобятся на последующем этапе моделирования исполнения. Матрица допустимости определяется как
\begin{equation}
 H_{ij}=\prod_{k=1}^{K}\mathbf 1\{c_{ik}\geq a_{jk}\}.
\end{equation}
Отсутствие требуемой возможности запрещает назначение независимо от ценности задачи. Количественные ограничения, например грузоподъёмность, при расширении модели должны задаваться отдельными неравенствами.

Маршрут $p_i=(p_{i1},\ldots,p_{i\ell_i})$ содержит различные идентификаторы задач. Его длина с возвратом в исходную точку равна
\begin{equation}
\begin{aligned}
 D_i(p_i)={}&\|s_i-r_{p_{i1}}\|
 +\sum_{u=1}^{\ell_i-1}\|r_{p_{iu}}-r_{p_{i,u+1}}\|\\
 &+\|r_{p_{i\ell_i}}-s_i\|,
\end{aligned}
\end{equation}
а для пустого маршрута принимается $D_i(\varnothing)=0$. Величина $E_i$ является ограничением длины пути, используемым как упрощённый ресурс; она не заменяет энергетическую модель БПЛА.

Пусть $x_{ij}=1$, если задача входит в маршрут аппарата $i$. Эталонная централизованная запись цели служит только для определения качества:
\begin{equation}
 \max_{p_1,\ldots,p_n}\ \sum_{i=1}^{n} S_i(p_i),\qquad
 S_i(p_i)=\sum_{j\in p_i} w_j-\lambda D_i(p_i),
\end{equation}
где $\lambda\geq0$ переводит длину маршрута в единицы полезности. Ограничения имеют вид
\begin{equation}
\begin{aligned}
 &x_{ij}\in\{0,1\},\quad x_{ij}\leq H_{ij},\quad
 \sum_i x_{ij}\leq1,\\
 &|p_i|\leq L_i,\qquad D_i(p_i)\leq E_i.
\end{aligned}
\end{equation}
Не все задачи обязательно получают исполнителя: этому могут препятствовать возможности аппаратов, ресурс или неположительная предельная полезность. Суммарная полезность и доля назначенных задач поэтому оцениваются раздельно.

\subsection{Сеть и доступная агенту информация}

Потенциальная связь задаётся неориентированным графом $G(t)=(\mathcal A,\mathcal E(t))$. Ребро возможно, если расстояние между аппаратами не превышает радиус $R$ и линия не заблокирована сценарием. Доставка в каждом направлении может независимо не состояться с вероятностью $p_{\rm loss}$. Таким образом, успешный обмен в отдельном такте может быть несимметричным.

Начальная модель использует дискретные синхронные такты. Пакет доставляется в текущем такте либо теряется; задержанные пакеты и их произвольное переупорядочивание пока не рассматриваются. Каждый агент знает список задач и собственные характеристики. Сведения о чужих планах доступны только из доставленных сообщений. Журнал симулятора используется для оценки метрик и не передаётся агентам.

В основной серии аппараты неподвижны, а сеть сначала разделена на компоненты, которые независимо формируют назначения. Затем связь восстанавливается. Такой режим позволяет исследовать восстановление согласованности без одновременного изменения стоимости маршрутов из-за движения аппаратов.

\section{Исследуемое правило перепланирования}

\subsection{Построение пакета и обязательный пересмотр}

Агент хранит пакет $b_i$, маршрут $p_i$, векторы предполагаемых победителей $z_i$, их ставок $y_i$ и сведения об актуальности информации по правилам CBBA~\cite{choi2009}. При построении пакета рассматривается предельная полезность вставки задачи:
\begin{equation}
 \Delta_{ij}=\max_{u\in\mathcal I_{ij}}
 \bigl[S_i(p_i\oplus_u j)-S_i(p_i)\bigr],
\end{equation}
где $\mathcal I_{ij}$ включает только вставки новой задачи, удовлетворяющие ограничениям. При пустом множестве вставок или $H_{ij}=0$ ставка недопустима. Из задач с положительной ставкой, выигрывающей у известного победителя, выбирается задача с максимальной предельной полезностью. При равенстве ставок выигрывает меньший идентификатор агента; при равных оценках собственных кандидатов выбирается меньший идентификатор задачи. Маршрут изменяется в месте вставки, а идентификатор задачи добавляется в конец пакета.

После согласования, если $k$ --- наименьшая позиция потерянной задачи в пакете, сохраняется только префикс
\begin{equation}
 b_i\leftarrow(b_{i1},\ldots,b_{i,k-1}).
\end{equation}
Из маршрута удаляются задачи освобождённого суффикса по их идентификаторам. Ограничивать такой сброс числом $q$ нельзя: более поздние ставки зависят от утраченного элемента. Полученная информация о победителях и отзыве собственных ставок распространяется с сохранением предусмотренных CBBA сведений об актуальности.

\subsection{Добровольный пересмотр после восстановления связи}

В начальном состоянии агенту известны только непосредственно доступные соседи; их сессии инициализируются без добровольного сброса. После старта появление новой линии или восстановление прежней обрабатывается одинаково. Сосед помечается недоступным после пяти тактов без сообщений; отдельная потеря пакета не завершает сессию. После двустороннего подтверждения связи формируется событие $e=(i,k,\nu_i,\nu_k)$, где $i<k$, а $\nu_i,\nu_k$ --- возрастающие номера сессий. Повтор подтверждения сохраняет те же номера. Сначала выполняются обмен состояниями и обязательное разрешение обнаруженных конфликтов.

Дополнительный пересмотр разрешён только двум концам восстановленной линии. Каждый из них однократно для события $e$ добровольно освобождает последние
\begin{equation}
 q_i(e)=\min\{q,|b_i|\}
\end{equation}
элементов оставшегося пакета, после чего снова выполняет построение пакета. Собственные ставки освобождённых задач отзываются; достоверные записи о чужих победителях не стираются. Обновлённые победители, ставки и сведения об их актуальности входят в очередной снимок состояния. Исследуются $q\in\{0,1,2\}$. Одновременные события одного агента объединяются в один добровольный сброс за такт; их идентификаторы запоминаются. Лимит действует на агента и такт: при подтверждении разных линий в разные такты суммарный пересмотр может быть больше $q$. Вариант $q=0$ совпадает с CBBA-контролем при общей сетевой обёртке.

Агенты за пределами этой пары продолжают обмен и согласование. Если до них дошёл конфликт, обязательное освобождение зависимого суффикса выполняется полностью. Поэтому локальность относится к запуску добровольного пересмотра, а не к гарантированной границе всех изменений в рое.

Исчезновение соседа не служит основанием считать его отказавшим или отменять его назначения. Периодическая передача состояния продолжается и после обработки события, чтобы потеря одного сообщения не прекращала согласование. Временные сведения об исходном владельце ставки нельзя заменять временем получения последнего пакета: поздняя пересылка старой информации не должна отменять более свежий отзыв.

Планируемый вклад состоит в исследовании этого правила при слиянии компонентов гетерогенной сети и в определении компромисса между объёмом пересмотра и качеством плана. Сходимость при бесконечных переключениях графа, гарантированное преимущество и глобальная оптимальность не утверждаются. Даже после прекращения переключений выбранная маршрутная полезность требует отдельной проверки свойств; в экспериментах учитываются незавершившиеся согласования.

\section{План экспериментов}

\subsection{Сценарии и сравниваемые варианты}

Планируется генерация синтетических сценариев в квадратной области размером $1000\times1000$~м. Задачи и аппараты размещаются равномерно. При $K=3$ требования задач равновероятно выбираются из векторов $(1,0,0)$, $(0,1,0)$, $(0,0,1)$. Сравниваются группа с возможностями $(1,1,1)$ у каждого аппарата и группы однофункциональных аппаратов с долями типов $1/3,1/3,1/3$ и $0{,}7;0{,}2;0{,}1$. Правило округления численности фиксируется в протоколе.

Радиус связи задаётся относительно $R_{\rm conn}$ --- минимального радиуса связности геометрического графа. Удаление самого длинного ребра его минимального остовного дерева задаёт две связные группы. Сначала все линии между ними блокируются. В основной серии восстанавливается только удалённая соединяющая линия, а остальные межгрупповые линии остаются заблокированными. В стрессовой серии одновременно открываются все такие линии: число инициаторов тогда может приближаться к размеру роя. Оно и суммарный добровольный сброс измеряются отдельно. Ещё один контроль сохраняет блокировку до конца запуска, без предположения о глобальном согласовании.

\begin{table}[htbp]
\caption{Начальные параметры основной серии}
\label{tab:protocol}
\centering
\small
\setlength{\tabcolsep}{3pt}
\begin{tabular}{p{0.44\columnwidth}p{0.46\columnwidth}}
\toprule
Параметр & Значения \\
\midrule
Число аппаратов $n$ & $10,20,40$ \\
Число задач $m$ & $2n,4n$ \\
Размер пакета $L_i$ & $5$ \\
Ресурс пути $E_i$ & $3000$ м \\
Радиус $R/R_{\rm conn}$ & $1{,}1;1{,}5$ \\
Вероятность потери & $0;0{,}1;0{,}3$ \\
Восстановление связи & такт $30$ или $90$ \\
Предел после восстановления & $300$ тактов \\
Число повторов & $30$ на конфигурацию \\
\bottomrule
\end{tabular}
\end{table}

Ценность задачи равна $100$, коэффициент $\lambda=0{,}01$ единицы полезности на метр. Эти числа являются параметрами будущего эксперимента. Перед итоговым сравнением их проверяют на отдельных сценариях настройки и затем фиксируют. Отсутствие допустимых назначений регистрируется, а не устраняется скрытой перегенерацией примера.

Сопоставляются обычное согласование CBBA, предлагаемое правило и адаптация частичного сброса CBBA-PR к восстановлению связи. В последней каждый получивший уведомление агент однократно освобождает до $q$ последних элементов пакета; одновременные события объединяются по тому же правилу. Доставка уведомлений учитывается в затратах. Это адаптация к сетевому событию, тогда как исходная работа рассматривает поступление задач. Дополнительная диагностическая абляция сопоставляет заранее заданные наборы участников при равном бюджете добровольного сброса; совпадение фактического числа освобождений проверяется по журналу.

В варианте Contract Net каждой задаче заранее сопоставляется агент-распорядитель по её идентификатору. Он собирает предложения и выдаёт контракт; единого распорядителя всех задач нет. Перед подтверждением исполнитель повторно проверяет допустимость актуального маршрута и может отклонить устаревшее предложение. Неназначенные задачи объявляются повторно. Изоляция распорядителя может задержать назначение и отражается в результатах. Для всех методов одинаковы требования задач, полезность, ресурсы и механизм потерь. Идеальная связь используется как отдельный ориентир.

\subsection{Метрики и воспроизводимость}

В статической серии измеряется доля задач с единственным допустимым исполнителем. Она отличается от доли реально выполненных задач. Число конфликтов определяется по текущим собственным планам агентов:
\begin{equation}
 C(t)=\sum_{j=1}^{m}\mathbf 1\left\{\sum_{i=1}^{n}x_{ij}(t)>1\right\}.
\end{equation}
Для семейства CBBA время восстановления оценивается ретроспективно: от снятия блокировки до начала последнего неизменного состояния, сохраняющегося до конца горизонта не менее десяти тактов. В нём все векторы победителей совпадают, соответствуют собственным планам и конфликты отсутствуют. Это конечный экспериментальный критерий, а не доказательство сходимости. Проверку выполняет наблюдатель; агенты её не используют. Contract Net не обязан хранить такие векторы, поэтому для него отдельно оценивается время подтверждения контрактов. Общие показатели всех методов сравниваются на фиксированном горизонте.

Дополнительно оцениваются суммарная длина маршрутов, полезность, число добровольных и обязательных освобождений, количество сообщений и переданные байты. В сетевые затраты входят служебные сообщения и повторные передачи, включая попытки, завершившиеся потерей. Затраты до восстановления согласованности и за фиксированный горизонт показываются отдельно. Полезность конфликтного плана не интерпретируется как качество допустимого решения.

Для каждого сценария сохраняются начальные данные, параметры, случайные начальные состояния генераторов и журнал обмена. Алгоритмы сравниваются попарно на одинаковых экземплярах. Потери индексируются тактом, направлением линии и номером попытки, чтобы различие числа сообщений не сдвигало общий случайный поток. Представляются медианы, межквартильные интервалы, доверительные интервалы парных разностей и доля запусков, достигших согласования. Незавершившиеся запуски не исключаются из отчёта.

Отдельный этап предусматривает движение аппаратов, обслуживание задач и оценку доли уникальных завершённых задач к заданному времени. Текущая выполняемая операция фиксируется, а будущие пакеты перестраиваются в общие для всех методов эпохи планирования: смена исходной позиции меняет стоимость последующих назначений. Информация о завершении становится доступна удалённому агенту только после доставки сообщения. Повторные исполнения регистрируются отдельно, поскольку последующее согласование планов не отменяет уже совершённые действия.

\section{Заключение}

Сформулирована задача назначения с учётом возможностей аппаратов, ограничений маршрута и временного разделения сети. Описан исследуемый вариант событийного частичного перепланирования: восстановление линии запускает добровольный пересмотр у её концов, а обязательное устранение конфликтов сохраняет общие правила работы с зависимыми ставками. Подготовлен протокол сравнения, позволяющий оценить качество распределения, время восстановления и стоимость обмена. Следующий этап включает реализацию методов и проведение экспериментов; по их результатам будут уточнены полезность предлагаемого правила и условия его применения.

\balance
\begin{thebibliography}{9}

\bibitem{choi2009}
Choi H.-L., Brunet L., How J.~P. Consensus-Based Decentralized Auctions for Robust Task Allocation // IEEE Transactions on Robotics. 2009. Vol.~25, no.~4. P.~912--926. DOI: \url{https://doi.org/10.1109/TRO.2009.2022423}.

\bibitem{buckman2019}
Buckman N., Choi H.-L., How J.~P. Partial Replanning for Decentralized Dynamic Task Allocation // AIAA Scitech 2019 Forum. 2019. AIAA~2019-0915. DOI: \url{https://doi.org/10.2514/6.2019-0915}. URL: \url{https://arxiv.org/abs/1806.04836}.

\bibitem{raja2022}
Raja S., Habibi G., How J.~P. Communication-Aware Consensus-Based Decentralized Task Allocation in Communication Constrained Environments // IEEE Access. 2022. Vol.~10. P.~19753--19767. DOI: \url{https://doi.org/10.1109/ACCESS.2021.3138857}.

\bibitem{smith1980}
Smith R.~G. The Contract Net Protocol: High-Level Communication and Control in a Distributed Problem Solver // IEEE Transactions on Computers. 1980. Vol.~C-29, no.~12. P.~1104--1113. DOI: \url{https://doi.org/10.1109/TC.1980.1675516}.

\bibitem{kha2025}
Kha A., Beynier A., Labreuche C., Marchand M. Enhancing CBBA convergence and optimality guarantees in Multiagent Task Allocation // ECAI~2025. 2025. P.~3485--3494. DOI: \url{https://doi.org/10.3233/FAIA251222}.

\bibitem{kiselev2025}
Киселев С.~Д., Алтунян А.~В. Локальные методы планирования траекторий на основе обучения с подкреплением // Искусственный интеллект в промышленных, коммерческих, медицинских и финансовых приложениях: сборник статей научно-технического семинара / под ред. А.~Р.~Ефимова. М.: НИТУ «МИСИС», 2025. Вып.~3. С.~84--89. URL: \url{https://misis.ru/files/-/7030ef7ab572d5227e50a55006446a76/kik_2024_2.pdf#page=85}.

\bibitem{bakhvalov2025}
Бахвалов И.~М., Старцев С.~В. Распознавание БПЛА различных классов средствами компьютерного зрения // Искусственный интеллект в промышленных, коммерческих, медицинских и финансовых приложениях: сборник статей научно-технического семинара / под ред. А.~Р.~Ефимова. М.: НИТУ «МИСИС», 2025. Вып.~3. С.~34--40. URL: \url{https://misis.ru/files/-/7030ef7ab572d5227e50a55006446a76/kik_2024_2.pdf#page=35}.

\bibitem{ly2024}
Ly K.~T., Semenov V., Risiglione M., Merkt W., Havoutis I. R-LGP: A Reachability-guided Logic-geometric Programming Framework for Optimal Task and Motion Planning on Mobile Manipulators // 2024 IEEE International Conference on Robotics and Automation (ICRA). 2024. P.~14917--14923. DOI: \url{https://doi.org/10.1109/ICRA57147.2024.10611389}. URL: \url{https://arxiv.org/abs/2310.02791}.

\end{thebibliography}
\end{document}
